\paragraph{Theorem D.2.49 (Proved: unbounded canonical star logs with the full integral boundary map).}
Let a marked zero-free star boundary have the integral basis
$(\alpha,\beta,\lambda,\ell)$ and monodromy
\[
\begin{gathered}
 A\alpha=-\alpha,\quad A\beta=-\beta-n\lambda,\quad
 A\lambda=\lambda-\alpha,\quad A\ell=\ell,\\
 \delta=-\alpha+2\lambda,\qquad\gamma^*(\ell)=1,
 \qquad n\in\{1,3\}.
\end{gathered}                                               \tag{D.SL1}
\]
There are twelve actual marked star positions of this type
in the natural catalogue: the eight primitive-gamma $I_1^*$
positions of D.2.27 and the four primitive-gamma $I_3^*$
positions. Their unimodular bases are retained in the
certificate below. The theorem also applies to any further
input for which the displayed integral marked basis is supplied.

Take a root of order $m=2k$, with $k>0$. Its rank-one
semistable Mumford cap has $q=nk$ components. Let its affine
deck action on their cyclic labels be
$i\mapsto j-i$, with a marked integer lift $j\in\mathbb Z$.
The component labels use $j$ modulo $q$; its integer lift is
retained in the peripheral marking. Write its collapsed
power as
\[
 v_{\rm col}=R\lambda+N\ell,\qquad R=2r-j,
                \qquad \gcd(N,2k)=1 .                     \tag{D.SL2}
\]
Every $j,r$ and every such $N$ is realized by an actual free
affine root action. All its holomorphic unit factors and
integer affine twists are retained. Put
\[
 f=\#\{i:2i=j\pmod q\},\qquad \rho=(q+f)/2.
\]
Thus $f=1$ for odd $k$, $f=2$ for even $k$ and even $j$,
and $f=0$ for even $k$ and odd $j$.
For the quotient cap $U$, the complete integral groups are
\[
\begin{array}{c|c}
a&H_a(U;\mathbb Z)\\\hline
0&\mathbb Z\\
1&\mathbb Z^2\\
2&\mathbb Z^{\rho+1}\oplus
                    (\mathbb Z/2\text{ if }f=0)\\
3&\mathbb Z^{2\rho}\\
4&\mathbb Z^\rho,
\end{array}                                                \tag{D.SL3}
\]
and higher groups vanish.

The boundary has $H_2(\partial U)=\mathbb Z^4$. Its
degree-two fibre coinvariant has the primitive basis
\[
 x=[\beta\lambda],\qquad
 w=[\beta\ell]+\frac{n-1}{2}[\lambda\ell],\qquad
 [\alpha\beta]=-2x,\quad [\beta\ell]=nw,\quad
 [\lambda\ell]=-2w.                                     \tag{D.SL4}
\]
Let $D,L$ be its literal fixed-normal tori in the
$\delta,\ell$ directions. In coordinates $(x,w,D,L)$,
the \emph{full} boundary kernel is as follows:
\[
\begin{array}{c|c}
f&\ker[H_2(\partial U)\longrightarrow H_2(U)]\\\hline
1&\{(0,b,c,d):kb=2Nc-Rd\}\\
2&\{(0,b,c,d):(k/2)b=Nc-(R/2)d\}\\
0&\{(0,b,c,d):kb=2Nc-Rd,\ b\equiv0\pmod2\}.
\end{array}                                                \tag{D.SL5}
\]
These are unsaturated marked kernels where appropriate.
Their projected indices in the full fixed lattice are
respectively $k,k/2,2k$. The boundary cokernel is free of
rank $\rho-1$. In particular the last row retains the
actual cap $\mathbb Z/2$ and its nonzero fibre correction;
it is not a conclusion from rational ranks or determinant
orders alone.

These columns can be inserted directly into the integral
bar algorithm of D.2.41 or D.2.46. A vector $(0,b,c,d)$
contributes the actual regular-fibre representative of
$bw$, together with $c$ times the marked $\delta$ sweep
and $d$ times the marked $\ell$ sweep. The cap contributes
$\rho$ to the free relative rank in (D.GI10) and zero to
the Euler number. Thus all other already proved cap
choices and full finite Thom columns can be retained.

There is also a closed all-order result when this is the
sole logarithmic cap and every other cap has the minimal
natural marking of D.2.26, with no extra regular logs.
Write the initial coefficient as $b_0=4B$, and let $e$
be the ordinary clutching integer. Put
\[
                         s=N-2ke .                       \tag{D.SL6}
\]
It is an odd nonzero integer, prime to $2k$. Then
\[
 \pi_1(X)=\mathbb Z/|s|,\qquad
 T=\operatorname{Tor}H_2(X)=
 \begin{cases}
 \mathbb Z/(2B|s|),&f=1,\\
 \mathbb Z/(B|s|),&f=2,\\
 \mathbb Z/2\oplus\mathbb Z/(2B|s|),&f=0.
 \end{cases}                                             \tag{D.SL7}
\]
The other groups are
\[
\begin{gathered}
 H_2=\mathbb Z^{\rho-1}\oplus T,\qquad
 H_3=\mathbb Z^{2\rho-2}\oplus T,\\
 H_1=\mathbb Z/|s|,\qquad
 H_4=\mathbb Z^{\rho-1}\oplus\mathbb Z/|s|,\qquad H_5=0,
 \qquad H_0=H_6=\mathbb Z .
\end{gathered}                                             \tag{D.SL8}
\]
Cyclic summands of order one are omitted. For example,
the six default natural rows with $b_0=4$ admit the
simply connected choice $s=\pm1$, $k=2h$, $j$ even.
It has
\[
          H_2=H_4=\mathbb Z^{nh},\qquad
                         H_3=\mathbb Z^{2nh}.              \tag{D.SL9}
\]
Here $n=1$ on OS50/default and $n=3$ on OS46/default.
These are nonminimal star-log outputs, and no smooth
classification is inferred from their homology.

\begin{proof}
\emph{The actual unbounded affine actions.}
From (D.SL1), $A^{2k}-I=q\alpha\beta^*$, so the
minimal semistable root model has exactly $q$ components.
This is the same integral Tate model used in D.2.24
and D.2.27, now at arbitrary root order. Its canonical
deck action reverses the component circle and fixes
the chosen zero component. Flat sections of the root
torus which are relevant to the affine action are
\[
                t=\frac{j}{q}\beta+\frac r k\lambda
                                      +\frac{N}{2k}\ell.
                                                               \tag{D.SL10}
\]
They are single-valued because
$(A^{2k}-I)t=j\alpha$ is integral. The first term
translates the Mumford fan by $j$ components and
therefore extends over the toric cap. The other
terms are ordinary unit translations in the
transverse periods. Their full norm is
\[
         \sum_{a=0}^{2k-1}A^a t
                     =-r\alpha+(2r-j)\lambda+N\ell.
                                                               \tag{D.SL11}
\]
It is integral, so composing the canonical deck
map with this translation still has order $2k$.
Using the inverse endpoint convention changes
only the coefficient of $\alpha$ in (D.SL11);
its collapsed vector (D.SL2) is unchanged, and
the corresponding exterior differences vanish
in the fibre coinvariant (D.SL4).

The gamma displacement of a nonidentity deck
power is nonintegral because $N$ is prime to
$2k$. Hence the action is free on every stratum,
including fixed components and fixed double
curves. Conversely, in this canonical primitive-
gamma sector the flat translation conditions
give $q t_\beta\in\mathbb Z$, $k t_\lambda\in\mathbb Z$
and $2k t_\ell\in\mathbb Z$; a remaining angular
translation is removed by conjugation since
$A\alpha=-\alpha$. This is the marking (D.SL10).
The usual root-and-unit isotopy keeps these
integral characters and all component meridians.
Thus the following topological calculation
applies to each actual action, not merely to
an abstract permutation of components.

\emph{The normalization and conductor at every order.}
The cap retracts equivariantly to its central
rank-one toric fibre. The quotient of its
$q$-gon dual graph by the reflection is an
interval. There are $\rho$ vertex orbits and
$\epsilon$ edge orbits, with
$\rho+\epsilon=q+1$; the number of fixed vertices
plus fixed edges is two. The normalization
and double-curve incidence graph is the
barycentric subdivision of this interval,
with exactly $q$ conductor incidences.

A paired vertex has stabilizer generated by
the square of the deck map. Its open
$\mathbb C^*$ bundle has an abelian period
lattice. The angular circle is primitive
because $\gcd(N,k)=1$. Thus its normalization
is a ruled surface over a torus with zero
Chern class, and its two pole sections have
the same class $\sigma$.
A fixed vertex instead has a pole-interchanging
clutch. Around the transverse $\lambda$ circle
the relation $A\lambda=\lambda-\alpha$ gives
the projective loop
$z\mapsto e^{-2\pi i u}/z$.
Its determinant winding is odd. Passing to
the primitive quotient gamma loop uses an
odd power of this map, since $N$ is odd, and
does not change the parity. Its normalization
is therefore the odd oriented $S^2$ bundle
over a torus. Choose a section $\sigma$ with
$\sigma^2=1$ and denote its complex sphere
fibre by $F$. Its connected conductor
bisection has class $2\sigma-F$; its
self-intersection is zero because the pole
normal line becomes trivial on the cover.

Every conductor flag has the even subgroup
as stabilizer. Its map to a normalization
base has degree two at a fixed vertex and
one at a paired vertex. Its map to the
double curve has degree two at a fixed edge
and one at a paired edge. Consequently
$H_2(U)$ has the following \emph{integral}
presentation. For each vertex orbit use
$F_v,\sigma_v$, and for each edge orbit use
its complex double-curve class $D_e$.
For each conductor incidence impose
\[
 \begin{cases}
 2\sigma_v-F_v-d_eD_e=0,&v\text{ fixed},\\
 \sigma_v-d_eD_e=0,&v\text{ paired},
 \end{cases}
 \qquad d_e=\begin{cases}2,&e\text{ fixed},\\1,&e\text{ paired}.
 \end{cases}                                             \tag{D.SL12}
\]
There is no hidden degree-one bridge in this
normalization sequence: the conductor map
on $H_1$ is injective, by leaf elimination
on the incidence tree, since every local
isogeny is injective on $H_1$. The map on
$H_2$ is likewise injective. Hence
$H_3(U)$ is the direct sum of the
$\mathbb Z^2$ groups of the normalizations,
and $H_4(U)$ is their direct sum of top
classes.

All conductor section classes in (D.SL12)
are a single common class $S$. Internal
paired vertices have $\sigma_v=S$;
fixed vertices satisfy $2\sigma_v-F_v=S$;
a fixed endpoint edge satisfies $2D_e=S$.
Leaf elimination now gives free rank
$\rho+1$. If a vertex is fixed its sphere
coefficient is a unit and no torsion
survives. If no vertex is fixed, the two
fixed endpoint edges give
\[
 S=2D_0=2D_1,\qquad
       \tau=D_1-D_0\text{ of exact order two}.             \tag{D.SL13}
\]
This proves the degree-two part of (D.SL3).

For clarity, the fundamental group of the
cap follows directly from the free-cover
description after collapsing $\alpha$:
\[
\begin{gathered}
 \Gamma=\langle\beta,\lambda,\ell,g\mid
 [\beta,\lambda]=[\beta,\ell]=[\lambda,\ell]=1,\\
 g\beta g^{-1}=-\beta-n\lambda,\quad
 [g,\lambda]=[g,\ell]=1,\quad
                    2kg=R\lambda+N\ell\rangle .
\end{gathered}                                             \tag{D.SL14}
\]
Additive notation in the last relation is
within the abelian subgroup generated by
$g,\lambda,\ell$. Its abelianization is free
of rank two: the relation $2\beta+n\lambda=0$
is primitive because $n$ is odd, and the
remaining relation is primitive because
$N$ is prime to $2k$. This proves $H_1(U)$
and finishes the local groups. As an
independent calculation, choose a primitive
quotient gamma circle. The cap is a mapping
torus of $I_q\times S^1_\lambda$; on its
component spheres the return is the
reflection, while on the actual specialized
torus $t=[\beta\lambda]$ naturality of the
exterior-square action gives
$t\mapsto-t-\sum_iF_i$. Its integral Wang
sequence gives the same groups, including
the order-two summand in (D.SL13).

\emph{The regular-fibre map and its missing parity.}
The coinvariant relations (D.SL4) follow
by direct integral exterior-square reduction
of (D.SL1); all four nonzero Smith factors
are units. The regular torus $\alpha\beta$
specializes to the sum of the $q$ complex
component spheres, each with degree one.
The transverse torus $\lambda\ell$ maps
to $kS$. Indeed its chosen component has
stabilizer of order $k$ in the paired case,
so its base cover has degree $k$ and its
class is $k\sigma$. In the fixed case its
base cover has degree $2k$; it is disjoint
from the bisection and has class
$k(2\sigma-F)$. Both give $kS$ in (D.SL12).
Thus
\[
        2x=-\sum_v|v|F_v,\qquad 2w=-kS.                  \tag{D.SL15}
\]
When $f>0$ the cap group is free, so these
equations determine the images uniquely.
Put
$T_0=\sum_{v\,\mathrm{paired}}F_v+
       \sum_{v\,\mathrm{fixed}}\sigma_v$.
For $f=1$ let $D_0$ be the sole fixed-edge
class. For $f=2$ retain the common $S$.
Then the images are
\[
\begin{array}{c|c|c}
 f&x&w\\\hline
1&-T_0+D_0&-kD_0\\
2&-T_0+S&-(k/2)S.
\end{array}                                               \tag{D.SL16}
\]
For $f=0$, equations (D.SL15) alone leave
two genuine order-two ambiguities. We
resolve them using actual conductor loops.

Choose integers $a_0,b_0$ with
$a_0N+b_0(2k)=1$ and put
$\kappa=a_0g+b_0\ell$. The classifying
group (D.SL14) becomes the mapping-torus
group of the two-torus $(\beta,\lambda)$
with return
\[
           B=\begin{pmatrix}-1&0\\-n&1\end{pmatrix}.
\]
In fact
$\ell=2k\kappa-a_0R\lambda$ and
$g=N\kappa+b_0R\lambda$.
Here $a_0$ is odd. Its integral degree-two
group is
\[
 H_2(\Gamma)=\mathbb Zz\oplus\mathbb Z/2\,t,
       \qquad z=[\lambda\kappa],\quad t=[\beta\lambda].
                                                               \tag{D.SL17}
\]
This group is computed by the same integral
bar comparison as D.2.41. For the two-dimensional
return $B$, its quadratic correction is
$f_B(u)=-n\binom{u_\beta}{2}t$.
Thus a $2k$-fold $\beta$ norm torus has
free part $-knz$ and torsion part $-knt$.
The actual regular tori therefore map as
\[
 [\beta\ell]\longmapsto-knz-(kn+a_0R)t,
 \qquad [\lambda\ell]\longmapsto2kz.                     \tag{D.SL18}
\]

First choose $0\le j<q$ and label the two fixed endpoint edges so that
their stabilizer loops are $g$ and $\beta g$.
These are the actual two reflected edge
lifts, separated by half the dual circle.
Their primitive transverse periods are
$\kappa$ and
$\beta\kappa-\frac{n(a_0-1)}2\lambda$, respectively.
Their complex torus classes consequently
map to
\[
                        D_0\mapsto z,\qquad D_1\mapsto z-t.
                                                               \tag{D.SL19}
\]
Both use the positive quotient gamma period,
so their orientations agree. It follows
that $\tau$ maps to the nonzero $t$.
After quotienting $H_2(U)$ by its component
spheres, (D.SL19) is an isomorphism onto
(D.SL17). In particular it detects precisely
the torsion ambiguity left in (D.SL15).

Since $k$ is even and $j,R,a_0$ are odd
in this case, (D.SL18) sends the primitive
$w$ of (D.SL4) to $-kz+t$; it sends $x$
to $t$. Hence the exact, marked images are
\[
       x=-\sum_vF_v+\tau,\qquad
                         w=-kD_0+\tau.                   \tag{D.SL20}
\]
Both order-two terms are forced by the
actual classifying map. In particular the
fibre cokernel has cyclic torsion
$\mathbb Z/(2k)$ in this case; replacing
the second image by $-kD_0$ would give a
different marked map.

\emph{The literal fixed-normal tori.}
The root meridian is $g^{2k}v^{-1}$,
with (D.SL11). Sweeping its bounding
normal disc by $\delta$ and $\ell$ gives
the exact relations in $H_2(U)$
\[
                 2kD+4Nw=0,\qquad
                             2kL-2Rw=0.                  \tag{D.SL21}
\]
The omitted $\alpha$ terms vanish by
(D.SL4). These relations determine the
free parts of both images. When $f>0$
there is no torsion ambiguity. When $f=0$,
the classifying map fixes it as well:
the cap loops are $\delta=2\lambda$,
$g=N\kappa+b_0R\lambda$ and
$\ell=2k\kappa-a_0R\lambda$, so their
actual commuting tori map to
\[
                 D\mapsto2Nz,\qquad L\mapsto-Rz
                            \quad\text{in }H_2(\Gamma).
\]
There is no $t$ term. In the same notation
as (D.SL16)--(D.SL20), the full images are
\[
\begin{array}{c|c|c}
f&D&L\\\hline
1&2ND_0&-RD_0\\
2&NS&-(R/2)S\\
0&2ND_0&-RD_0.
\end{array}                                               \tag{D.SL22}
\]
Together with the fibre images these
prove (D.SL5) by solving an integral
linear relation. For $f=0$ the coefficient
of $x$ must vanish in the independent
component-sphere direction; the coefficient
of $w$ must then be even by (D.SL20).
The free coefficient gives the displayed
equation. This explicitly retains the
relative torsion information rather than
saturating the normal lattice.

The cokernel of the full boundary map is
free of rank $\rho-1$, as one sees directly
from (D.SL12) and (D.SL16)--(D.SL22).
Equivalently, Poincare--Lefschetz duality
puts it into the free group $H^4(U)$.
The projected-kernel indices are $k,k/2,2k$;
the computations for $m=2$ and the swapped
$m=4$ in D.2.27 are recovered exactly.
The full local kernel can therefore be
used in the global bar presentation,
with all its existing endpoint corrections.

The integer lift of $j$ must also be retained.
Write $j=j_0+hq$, $0\le j_0<q$. The corresponding
peripheral generator changes from $g_0$ to
$g=\beta^h g_0$, and $R=R_0-hq$.
Its literal normal tori satisfy
\[
                    D=D_0^{\rm norm},\qquad
                    L=L_0-hn w.                           \tag{D.SL24}
\]
Here the first symbol denotes the old normal
torus, not a double-curve class. Indeed
$[\delta\beta]=0$ and $[\beta\ell]=nw$
in (D.SL4). Substituting (D.SL24) into
the kernel formulas replaces $R_0$ by
$R$ exactly. In the $f=0$ case it changes
the displayed image of $L$ in (D.SL22) to
$-RD_0+h\tau$, but the equation in
(D.SL5) already forces its coefficient
$d$ to be even; the kernel parity is
therefore still precisely $b$ even.
Thus integer affine lifts are included
without suppressing their actual torsion
correction.

\emph{The closed one-star Smith calculation.}
For the natural minimal exterior, the
fibre functional has
$\Omega(x)=-a$ and $\Omega(w)=b_0/4=B$.
The primitive graph calculation has the
single diagonal relation $D=-b_0e\,s_0$,
where $s_0$ is its regular-fibre generator.
Because every vector in (D.SL5) has zero
$x$ coordinate, its global relation on
$(s_0,L)$ is
\[
                         (B(b-4ec),d).                   \tag{D.SL23}
\]
Set $u=b-4ec$ and $s=N-2ke$.
For $f=1$ its defining condition becomes
$ku+Rd=2sc$. Since $k$ is odd and prime
to $2s$, $d$ is unrestricted and
$u$ is a fixed residue modulo $2s$.
The Smith cokernel is cyclic of order
$2B|s|$.
For $f=2$, the condition is
$(k/2)u+(R/2)d=sc$.
Here $k/2$ is prime to $s$, so the same
argument gives a cyclic cokernel of
order $B|s|$.
For $f=0$, $u$ and $d$ are even; write
$u=2u_0,d=2d_0$. The condition becomes
$ku_0+Rd_0=sc$. Its relation matrix has
Smith factors $2$ and $2B|s|$.
This proves the three torsion groups
in (D.SL7) for every parameter.

The one-log Seifert presentation gives
$\pi_1=\mathbb Z/|s|$. Its rational
Euler number cannot vanish because
$N$ is prime to $2k$. The free local
boundary defect is $\rho-1$ and there
is no additional rational degree-one
bridge. All other caps in the stated
corollary have their minimal ranks.
Thus $b_2=\rho-1$. The unique primitive
rank-two zero still contributes Euler
number two, while this star cap has
Euler number zero. Integral duality
gives (D.SL8), and the simply connected
specialization gives (D.SL9).
\end{proof}

The exact script \path{code/manifolds_star_log_caps.py}
retains the twelve actual canonical bases,
the full normalization matrices, the
literal fibre and normal maps, and their
complete integral kernels. It checks
the conductor calculation independently
against the mapping-torus sequence and
compares the global Smith formula on
unbounded symbolic lattices and a finite
set of exact instances. The proof above,
including the order-two marking, is
independent of the finite parameter
bounds used in those corroborating runs.
Root independently read the full proof and replayed all
$12$ actual marked bases, $95$ normalization comparisons
and $1638$ global Smith calculations successfully.
An independent adversarial review by the fourfolds
agent also checked (D.SL17)--(D.SL24), including the
nonzero order-two correction and the integer phase lift.
Root actions with nonprimitive gamma
power, other integral star blocks and
positive-weight star degenerations are
outside this theorem's stated scope.
